\chapter{Résultats et constats}
\label{chap:resultats}

Les constats ci-dessous sont uniquement ceux établis par l'EDA et confirmés par les visualisations revues. Ils décrivent le jeu étudié.

\begin{keyfindings}
\begin{itemize}
\item Les volumes sont globalement équilibrés entre entreprises, types de campagnes, audiences, localisations et segments (écarts de parts inférieurs à 0,4 point).
\item Le rythme est d'environ 548 campagnes par jour civil.
\item Février est plus faible en volume brut en raison de ses 28 jours, pas d'une baisse d'activité journalière distincte.
\item Le CTR moyen vaut 14,04~\% contre une médiane de 9,98~\% (asymétrie à droite).
\item Le ROI moyen est $\approx$ 5 (5,0024 ; médiane 5,0100) dans les différentes dimensions analysées.
\item Les écarts de ROI moyen entre catégories restent faibles (0,005 à 0,038 selon la dimension).
\item Les corrélations CTR/clicks ($r = +0,51$) et CTR/impressions ($r = -0,66$) sont liées à la définition du CTR.
\item Il n'apparaît pas de gradient clair du ROI selon le CTR, la durée ou l'engagement.
\end{itemize}
\end{keyfindings}

\begin{warningbox}[Pas de causalité]
Ces résultats sont descriptifs. Une corrélation, un écart de moyennes ou une différence de volume n'établit pas de mécanisme causal.
\end{warningbox}
